\documentclass[11pt,a4paper]{article}

\usepackage{fontspec}
\usepackage[french]{babel}
\usepackage{microtype}
\usepackage[margin=2.25cm]{geometry}
\usepackage{amsmath,amssymb}
\usepackage{booktabs,longtable,tabularx}
\usepackage{enumitem}
\usepackage{xcolor}
\usepackage{hyperref}
\usepackage{url}
\usepackage{fancyhdr}
\usepackage{titlesec}

\setmainfont{Latin Modern Roman}
\setsansfont{Latin Modern Sans}
\setmonofont{Latin Modern Mono}

\definecolor{ink}{HTML}{18212B}
\definecolor{navy}{HTML}{173B57}
\definecolor{blue}{HTML}{176B91}
\definecolor{rulegray}{HTML}{CFD8DE}

\color{ink}
\hypersetup{
  colorlinks=true,
  linkcolor=navy,
  urlcolor=blue,
  pdftitle={Prompt Fable 5 — Étude de sensibilité du modèle M4B},
  pdfauthor={Anatole Joseph Stouls}
}
\urlstyle{tt}
\setlist[itemize]{leftmargin=1.5em,itemsep=0.3em,topsep=0.4em}
\setlist[enumerate]{leftmargin=1.7em,itemsep=0.35em,topsep=0.4em}
\setlength{\parindent}{0pt}
\setlength{\parskip}{0.55em}
\setlength{\emergencystretch}{3em}

\titleformat{\section}
  {\Large\bfseries\sffamily\color{navy}}{\thesection}{0.65em}{}
\titleformat{\subsection}
  {\large\bfseries\sffamily\color{blue}}{\thesubsection}{0.55em}{}
\titlespacing*{\section}{0pt}{1.7em}{0.55em}
\titlespacing*{\subsection}{0pt}{1.2em}{0.4em}

\pagestyle{fancy}
\setlength{\headheight}{22pt}
\fancyhf{}
\fancyhead[L]{\small\sffamily Prompt Fable 5}
\fancyhead[R]{\small\sffamily Sensibilité M4B}
\fancyfoot[C]{\small\sffamily\thepage}
\renewcommand{\headrulewidth}{0.4pt}
\renewcommand{\headrule}{\hbox to\headwidth{\color{rulegray}\leaders\hrule height \headrulewidth\hfill}}

\newcommand{\code}[1]{\texttt{\color{navy}#1}}
\newcommand{\file}[1]{\nolinkurl{#1}}
\newcommand{\term}[1]{\textbf{\color{navy}#1}}
\newenvironment{briefbox}
  {\begin{center}\begin{minipage}{0.94\textwidth}\color{navy}\hrule\vspace{0.7em}}
  {\vspace{0.5em}\hrule\end{minipage}\end{center}}
\newenvironment{notebox}
  {\begin{center}\begin{minipage}{0.90\textwidth}\color{navy}\itshape}
  {\end{minipage}\end{center}}

\begin{document}

\begin{titlepage}
  \centering
  \vspace*{2.2cm}
  {\sffamily\bfseries\color{navy}\fontsize{28}{34}\selectfont
    Étude de sensibilité\par}
  \vspace{0.35cm}
  {\sffamily\bfseries\color{blue}\fontsize{24}{30}\selectfont
    du modèle M4B\par}
  \vspace{1.1cm}
  {\Large Prompt de recherche pour Claude Fable 5\par}
  \vspace{1.6cm}
  \begin{minipage}{0.78\textwidth}
    \begin{briefbox}
      Mission de recherche numérique de bout en bout : plan d'expérience,
      simulations stochastiques, analyse de sensibilité, contrôle de robustesse
      et rapports scientifiques reproductibles. Le modèle étudié est M4B.
    \end{briefbox}
  \end{minipage}
  \vfill
  {\large Anatole Joseph Stouls\par}
  \vspace{0.25cm}
  {\large 17 juillet 2026\par}
  \vspace{0.8cm}
  {\small\sffamily Version révisée après annotations du lecteur.\par}
  {\small\sffamily À exécuter de préférence avec l'effort \code{xhigh}.\par}
\end{titlepage}

\tableofcontents
\clearpage

\section{Rôle et objectif scientifique}

Tu es responsable d'une étude de sensibilité reproductible du modèle de
société de crédit \term{M4B}. Travaille comme chercheur en économophysique,
systèmes complexes et simulation stochastique : audite le moteur, fige un
protocole, exécute la campagne, analyse les résultats et rédige les rapports
scientifiques en français.

L'étude doit cartographier honnêtement :

\begin{enumerate}
  \item quels paramètres contrôlent le régime démographique, le réseau de
    crédit, les faillites en cascade et les distributions individuelles ;
  \item où se trouvent les non-linéarités, seuils, interactions et changements
    de régime ;
  \item quelles signatures sont robustes aux paramètres, aux graines, à la
    durée et aux choix de fenêtre statistique ;
  \item quels paramètres semblent actifs, redondants ou assimilables à des
    conventions d'échelle, afin d'éclairer une éventuelle réduction future du
    modèle.
\end{enumerate}

L'objectif n'est ni de chercher la configuration qui « donne le plus de SOC »,
ni de sélectionner les seuls résultats compatibles avec le rapport antérieur.
Les extinctions, explosions, absences d'effet et contradictions sont des
résultats à conserver et à expliquer.

\begin{notebox}
\textbf{Consigne d'autonomie.} Quand tu as assez d'information pour agir,
agis. Mène la mission jusqu'aux simulations, aux analyses et aux rapports
compilés ; ne termine pas sur un simple plan ou une promesse d'exécution.
\end{notebox}

\section{Périmètre : M4B uniquement}

Le modèle étudié est :

\begin{center}
  \file{/home/anatole/jupyter/m4b_credit_soc_mini/}
\end{center}

M4B n'est pas traité comme un simple auxiliaire de M4 : il constitue l'objet
scientifique de cette campagne. Les simulations de sensibilité, les cartes de
régime, les conclusions et les recommandations de réduction doivent être
obtenues depuis M4B.

\subsection{Recours exceptionnel à M4}

Le dossier \file{m4_credit_soc_fable/} est une référence historique et de
provenance, pas un second modèle à balayer. Il n'y a \term{aucune campagne M4}
à mener en parallèle.

Tout nouveau recours exécutable à M4 doit satisfaire les quatre conditions
suivantes :

\begin{enumerate}
  \item une question précise ne peut pas être résolue à partir du moteur M4B,
    de ses sorties, de ses tests ou de son rapport mécanique ;
  \item le besoin et la conclusion attendue sont écrits avant l'exécution ;
  \item le nombre de runs M4 est minimal et leur résultat est isolé dans une
    annexe méthodologique, sans entrer dans les estimateurs M4B ;
  \item le rapport explique pourquoi ce recours était indispensable et ce
    qu'il permet de conclure sur M4B.
\end{enumerate}

Le test de parité existant constitue un recours justifiable une fois comme
contrôle de provenance de la réduction. Il ne transforme pas M4 en objet de
l'étude et ne doit pas être répété à chaque cellule. Si la preuve de parité
existante et l'intégrité du code suffisent, cite-les sans relancer M4.

La lecture d'un rapport M4 pour comprendre l'histoire du mécanisme est permise,
mais les affirmations de la présente étude doivent être vérifiées sur M4B.

\section{Répertoire de travail et lectures obligatoires}

Travaille dans \file{/home/anatole/jupyter}. Utilise le Python du venv :
\file{/home/anatole/jupyter/.venv/bin/python3}.

Lis intégralement avant de figer le protocole :

\begin{itemize}
  \item guide Fable :
    \file{/home/anatole/jupyter/m4_credit_soc/Prompting Claude Fable.md} ;
  \item présentation et usage de M4B :
    \file{m4b_credit_soc_mini/README.md} ;
  \item rapport mécanique M4B :
    \file{m4b_credit_soc_mini/report/mecanique_m4b.pdf} et sa source
    \code{mecanique\_m4b.tex} ;
  \item moteur et formats de sortie : \file{m4b_credit_soc_mini/m4b/},
    \file{run.py} et \file{tests/} ;
  \item intégration de M4B dans Simulation Lab :
    \file{modeles-systeme-physicoeconomique/m4b_credit_soc_mini/} et
    \file{simulation_lab/}.
\end{itemize}

Crée tous les scripts, données dérivées, figures, notes techniques et rapports
de l'étude dans :

\begin{center}
  \file{/home/anatole/jupyter/recherche/sensibilite_m4b/}
\end{center}

Traite comme \term{source en lecture seule} le moteur
\file{m4b_credit_soc_mini/m4b/}, ses tests, son rapport et les résultats
antérieurs. Une lacune d'instrumentation doit être contournée dans les scripts
de campagne ou signalée ; elle ne justifie pas silencieusement une nouvelle
version de la dynamique.

\section{Paramètres et statut expérimental}

Les paramètres scientifiques publics de M4B sont :

\begin{longtable}{@{}p{0.18\textwidth}p{0.70\textwidth}@{}}
\toprule
\textbf{Paramètre} & \textbf{Rôle} \\
\midrule
\endhead
\code{lam} & Intensité des naissances et axe de taille démographique. \\
\code{delta} & Taux de dépréciation du capital. \\
\code{sigma} & Volatilité du choc multiplicatif. \\
\code{K0} & Capital attribué à la naissance. \\
\code{k} & Taille de l'échantillon du marché local. \\
\bottomrule
\end{longtable}

\code{seed} définit une réplication stochastique. \code{T} définit la durée
d'observation. \code{pop\_max} est un garde-fou d'arrêt. Aucun des trois ne doit
être interprété comme un mécanisme économique. \code{alpha=1} et les
tolérances numériques sont des conventions ou constantes et ne doivent pas
être promues artificiellement en paramètres scientifiques.

\subsection{Rôle strict de l'étude temporelle}

La variation de \code{T}, du burn-in et des fenêtres sert à vérifier la
\term{probité des résultats} : convergence, stationnarité, stabilité des
estimateurs et absence de conclusion fondée sur un transitoire. Elle ne doit
pas devenir un axe de sensibilité économique ni entrer dans le classement des
paramètres scientifiques.

Une conclusion n'est robuste dans le temps que si :

\begin{itemize}
  \item elle persiste lorsque l'horizon est prolongé ;
  \item elle ne dépend pas d'une seule définition du burn-in ;
  \item les métriques cumulatives et les métriques par fenêtres donnent une
    lecture compatible ;
  \item les dérives résiduelles sont quantifiées et visibles.
\end{itemize}

\section{Grandeurs de réponse}

Avant les runs de production, définis dans le rapport de protocole un
dictionnaire des métriques : formule, unité, fenêtre temporelle, condition de
validité et statut primaire ou secondaire.

\subsection{Régime et démographie}

\begin{itemize}
  \item statut \code{ok}, extinction ou arrêt par \code{pop\_max} ;
  \item population moyenne après burn-in, $N/\lambda$, pente temporelle,
    autocorrélation et variabilité ;
  \item naissances, morts, durée de vie, renouvellement et bilan des flux ;
  \item stationnarité vérifiée sur plusieurs fenêtres et plusieurs horizons,
    jamais sur la seule valeur finale.
\end{itemize}

\subsection{Crédit, exposition et réseau}

\begin{itemize}
  \item nombre et volume de prêts, dette rapportée au capital ou aux actifs,
    degrés et concentration des expositions ;
  \item rythme d'accumulation puis d'élagage du carnet autour des grands
    événements ;
  \item métriques intensives ou normalisées par la population lorsque cela est
    requis par la comparaison.
\end{itemize}

\subsection{Avalanches et propagation causale}

\begin{itemize}
  \item fréquence, quantiles, maximum absolu et maximum rapporté à la
    population ;
  \item rapport de branchement, fraction induite, racines/taille et profondeur ;
  \item susceptibilité ou échelle de coupure, par exemple
    $\langle s^2\rangle/\langle s\rangle$ ;
  \item MLE \term{discret} de la loi de puissance, notamment au seuil
    $s_{\min}=2$, ajustement loi de puissance avec coupure et comparaison de
    vraisemblance avec une log-normale discrète renormalisée ;
  \item scaling de la coupure avec la taille du système.
\end{itemize}

Ne conclus jamais à une loi de puissance sur le seul $r^2$ d'une régression
log-log. Si un run contient trop peu d'avalanches pour un ajustement fiable,
marque la métrique non identifiable au lieu de produire un nombre trompeur.

\subsection{Distributions individuelles et inégalités}

\begin{itemize}
  \item capital $K$, valeur nette, revenu brut/net, Gini et renouvellement du
    haut de la distribution ;
  \item familles comparées sur les mêmes supports et avec les mêmes règles de
    troncature ;
  \item contrôle par âge ou cohorte lorsque l'interprétation
    distributionnelle peut en dépendre.
\end{itemize}

\subsection{Intégrité numérique}

\begin{itemize}
  \item égalité créances-dettes, absence de contrats orphelins, capital non
    négatif, bilan réel par pas et reproductibilité ;
  \item temps de calcul, volume des données et tout arrêt anticipé ;
  \item neutralité des options de mesure sur la trajectoire aléatoire.
\end{itemize}

\section{Protocole expérimental}

\subsection{Audit et réutilisation}

\begin{enumerate}
  \item Exécute les contrôles unitaires et comptables propres à M4B. Le test de
    parité M4 peut être cité ou lancé une fois selon les règles de la
    section~2.1.
  \item Inventorie les runs existants. Ne réutilise un run que si son moteur,
    sa configuration complète, sa graine, son horizon, son statut et ses
    données primaires sont vérifiables. Un nom de dossier ne suffit pas.
  \item Lance un benchmark représentatif pour mesurer coût et taille disque,
    puis fixe un budget défendable. Laisse au moins deux cœurs au système et
    évite la surallocation mémoire ou disque.
  \item Fige le protocole exploratoire et les règles de passage à la phase
    confirmatoire dans un rapport LaTeX daté.
\end{enumerate}

\subsection{Séparer taille, temps et paramètres microscopiques}

Traite \code{lam} d'abord comme axe de taille finie et démographique. Traite
\code{T} uniquement comme contrôle de probité temporelle. Ne les mélange pas
sans précaution dans un classement d'importance avec \code{delta},
\code{sigma}, \code{K0} et \code{k}.

La référence est :

\begin{center}
\code{lam=10, delta=0.05, sigma=0.25, K0=25, k=3, T=2000}.
\end{center}

Pour obtenir des queues mieux identifiées, le centre de la campagne peut être
\code{lam=30} ; justifie ce choix et conserve un lien explicite avec la
baseline.

Construis successivement :

\begin{enumerate}
  \item un contrôle de taille autour de $\lambda\in\{10,30,100\}$ ;
  \item un contrôle temporel limité à la vérification des résultats, par
    exemple $T\in\{2000,4000\}$ et plusieurs burn-ins ;
  \item des courbes OAT autour du centre pour rendre visibles monotonies,
    seuils et formes non linéaires ;
  \item un plan global espace-remplissant sur les paramètres microscopiques ;
  \item des coupes 2D ciblées choisies après le screening ;
  \item une phase confirmatoire indépendante sur les régimes et interactions
    les plus importants.
\end{enumerate}

\subsection{Plages initiales}

Ces plages doivent être auditées par des pilotes, puis corrigées si elles ne
produisent que des cas triviaux ou des arrêts :

\begin{center}
\begin{tabular}{@{}lll@{}}
\toprule
\textbf{Paramètre} & \textbf{Plage initiale} & \textbf{Remarque} \\
\midrule
\code{delta} & $0{,}02$--$0{,}10$ & autour de $0{,}05$ \\
\code{sigma} & $0{,}10$--$0{,}50$ & contrôle $\sigma=0$ séparé \\
\code{K0} & $5$--$100$ & échelle logarithmique préférable \\
\code{k} & $\{2,3,4,6,10\}$ & paramètre discret \\
\code{lam} & $\{10,30,100\}$ & scaling principal \\
\bottomrule
\end{tabular}
\end{center}

Ces bornes sont des points de départ, pas des vérités. Relie-les aux équations,
au rapport mécanique, aux pilotes et aux régimes observés.

Pour le plan global, utilise un plan reproductible adapté au mélange continu
et discret, par exemple un Latin hypercube stratifié pour les continus avec
équilibrage de \code{k}. Un ordre de grandeur de 48 à 80 points avec trois
graines par point est raisonnable après benchmark ; ajuste ce nombre selon le
coût et la précision observés. Ne revendique pas d'indices de Sobol si le plan
de Saltelli et le nombre de réplications ne les rendent pas valides.

Note pratique sur le temps de calcul : pour des valeurs petites de
\code{sigma} et \code{delta}, le système peut devenir très grand et mettre
beaucoup de temps à atteindre son régime permanent. Il peut alors être
nécessaire de pousser \code{T} bien au-delà de 10000, voire jusqu'à 30000.
Adapter \code{T} dans ce cas reste une exigence de probité temporelle, mais il
faut le faire avec prudence : attention au volume total des simulations et aux
coûts de stockage associés.

Utilise le même panel de graines pour chaque cellule d'une comparaison. Cela
permet des contrastes appariés, mais ne l'appelle pas abusivement « common
random numbers » si les trajectoires consomment ensuite des tirages différents.
Ne fusionne jamais les événements de plusieurs graines avant d'ajuster une
loi : ajuste chaque graine, puis agrège les estimateurs et leur incertitude.

Après screening, confirme les cellules retenues avec $T=4000$ et au moins cinq
graines. Près d'une frontière de régime, ou lorsque la variance inter-graines
domine l'effet, augmente les répétitions. Utilise par défaut un burn-in de
$T/4$ et vérifie la robustesse à la fenêtre.

\subsection{Diagnostic obligatoire du garde-fou de population}

\code{pop\_max} ne borne pas la dynamique : il censure une trajectoire par
sécurité. Atteindre cette valeur ne suffit donc jamais à conclure à une
explosion.

Pour toute cellule qui atteint \code{pop\_max} :

\begin{enumerate}
  \item répète-la avec plusieurs plafonds plus élevés et le même panel de
    graines ;
  \item examine les pentes, accélérations et temps de franchissement plutôt que
    la seule population finale ;
  \item encadre adaptativement la frontière dans l'espace des paramètres avec
    des cellules voisines des deux côtés ;
  \item vérifie qu'il existe un changement qualitatif entre un régime voisin
    stationnaire, qui peut frôler le plafond sans l'atteindre durablement, et
    un régime à croissance soutenue qui franchit successivement les plafonds ;
  \item si la moyenne stationnaire augmente continûment jusqu'à simplement
    croiser un plafond arbitraire, classe le run « censuré par le garde-fou » et
    non « explosif ».
\end{enumerate}

Le rapport doit montrer le diagnostic par des trajectoires et une carte locale,
et distinguer formellement \emph{changement de régime}, \emph{forte population
stationnaire} et \emph{censure numérique}.

\subsection{Analyse de sensibilité}

Produis plusieurs lectures complémentaires :

\begin{itemize}
  \item courbes OAT avec incertitude inter-graines et contrastes appariés ;
  \item corrélations de rang partielles ou régressions standardisées avec
    intervalles bootstrap pour le screening global ;
  \item modèle de réponse non linéaire validé hors échantillon pour détecter
    seuils et interactions, avec importance par permutation si pertinente ;
  \item cartes de régime extinction/stationnaire/croissance et diagrammes 2D
    des interactions confirmées ;
  \item décomposition explicite de la variabilité due aux paramètres et de la
    variabilité stochastique inter-graines.
\end{itemize}

Ne donne pas un classement unique sans préciser la métrique. Rapporte tailles
d'effet, incertitudes, non-monotonies et conditions de validité. Toute analyse
exploratoire utilisée pour choisir des cellules doit être étiquetée comme telle
et séparée de la confirmation.

\section{Infrastructure et traçabilité}

Organise le dossier de campagne au minimum ainsi :

\begin{verbatim}
recherche/sensibilite_m4b/
  scripts/
  manifests/
  results/
  figures/
  report/
    protocole.tex
    rapport_preliminaire.tex
    rapport_final.tex
\end{verbatim}

Exigences :

\begin{itemize}
  \item scripts relançables et reprise après interruption ;
  \item identifiant de run déterministe ou table de correspondance non ambiguë ;
  \item manifeste comportant version du moteur, tous les paramètres, graine,
    horizon, burn-in, statut, durée, chemin des données et contrôles
    d'intégrité ;
  \item résultats agrégés en CSV ou JSON pour les calculs, mais interprétation
    scientifique dans les rapports LaTeX ;
  \item journal technique des décisions, résultats négatifs et anomalies ;
  \item aucune cellule silencieusement exclue parce qu'elle s'éteint, atteint
    un garde-fou ou contredit l'hypothèse.
\end{itemize}

M4B est le modèle actif de Simulation Lab. Les runs confirmatoires retenus
doivent y être consultables, avec leurs figures et métadonnées, et être marqués
à conserver. Pour le screening massif, limite l'écriture individuelle
(\code{individual\_every=0} ou fréquence espacée) et les instantanés lorsque
ces données ne sont pas requises. Vérifie sur un contrôle que les options de
mesure ne changent pas la trajectoire.

\section{Livrables scientifiques en LaTeX}

Les livrables destinés au lecteur sont des \term{rapports écrits en LaTeX et
compilés en PDF}. Les fichiers CSV, JSON, scripts et manifestes sont des annexes
techniques, pas des substituts au raisonnement écrit.

Rédige au minimum :

\begin{enumerate}
  \item \code{protocole.tex/.pdf} : modèle, questions, paramètres, métriques,
    plan d'expérience, règles de confirmation, diagnostic de \code{pop\_max}
    et critères d'arrêt ;
  \item \code{rapport\_preliminaire.tex/.pdf} : audit, pilotes, qualité des
    données, bornes révisées, puissance statistique et décisions avant la
    confirmation ;
  \item \code{rapport\_final.tex/.pdf} : méthode, résultats, résultats
    négatifs, cartes de régime, incertitudes, limites, recommandations de
    réduction et perspective M5.
\end{enumerate}

Ces documents s'adressent à un \term{public avisé} en économie quantitative,
économophysique ou systèmes complexes. Ils doivent rester autonomes et
compréhensibles sans lire le code ni les échanges de l'agent :

\begin{itemize}
  \item présenter les équations utiles et la chronologie de M4B ;
  \item définir toute métrique avant son interprétation ;
  \item distinguer faits simulés, inférences statistiques et hypothèses ;
  \item expliquer les choix de méthode et les limites sans jargon d'outillage ;
  \item commenter chaque tableau et chaque figure dans le texte ;
  \item donner les effectifs, graines, horizons, intervalles d'incertitude et
    provenance des données ;
  \item ne jamais renvoyer le lecteur à un CSV brut pour comprendre une
    conclusion.
\end{itemize}

Chaque figure et tableau doit être reproductible depuis le manifeste. Compile
les documents avec XeLaTeX ou LuaLaTeX, vérifie les références croisées, les
légendes, les débordements et la lisibilité des pages finales.

\section{Suite du projet : horizon M5, hors périmètre}

Le rapport final doit informer le lecteur de la suite envisagée, sans la mener
dans la présente étude.

\subsection{Réduction éventuelle avant M5}

La sensibilité M4B doit identifier les paramètres actifs, faiblement
identifiables, redondants ou essentiellement liés aux unités. Elle peut
recommander une réduction argumentée du nombre de paramètres. Elle ne doit pas
implémenter silencieusement cette réduction dans M4B : la décision appartient
à la conception de M5.

\subsection{Nouvel exposant de concavité}

Dans M4B, la production ou extraction est actuellement de type racine carrée :

\[
  F(K)=\alpha K^{1/2}.
\]

M5 introduira un exposant continu $\gamma$ :

\[
  F_{\gamma}(K)=\alpha K^{\gamma},
  \qquad 0<\gamma<1,
\]

avec $\gamma=1/2$ comme cas M4B et, par exemple, $\gamma=1/1{,}5=2/3$ ou
$\gamma=1/3$. Ce paramètre n'appartient pas à la campagne actuelle. Le rapport
M4B doit toutefois fournir les baselines, unités, régimes et métriques qui
permettront d'isoler son effet dans M5.

\subsection{Variation dynamique et effet rebond}

Une étape ultérieure cherchera à faire apparaître un \term{effet rebond} à
partir d'une variation dynamique des paramètres du système. Au sens économique
de Sorrell et Dimitropoulos (2008), une amélioration d'efficacité réduit le coût
effectif d'un service, stimule sa consommation et reprend une partie des
économies de ressource attendues ; si cette reprise dépasse la totalité des
économies prévues, la littérature parle de \emph{backfire}.

Référence locale :

\begin{center}
\file{/home/anatole/Zotero/storage/TMKGGCJA/}\par
\file{Sorrell et Dimitropoulos - 2008 - The rebound effect Microeconomic definitions, lim.pdf}
\end{center}

Dans M5, ce concept devra être traduit avec prudence : définir une amélioration
d'efficacité, la ressource consommée, le service ou l'activité utile, puis
comparer la réponse dynamique à un contrefactuel statique « sans rebond ». Une
simple hausse de population, de production ou de crédit après un changement de
paramètre ne suffira pas à établir un effet rebond.

La présente étude ne fait varier que des paramètres constants entre runs. Elle
ne doit donc ni simuler cette politique dynamique, ni annoncer un effet rebond.
Elle doit seulement livrer une carte de référence et signaler quelles variables
seront nécessaires à sa mesure future.

\subsection{Étude théorique du comportement du système}

M5 sera accompagné d'une étude analytique destinée à préciser le comportement
attendu avant ou en parallèle des simulations. Elle devra notamment examiner :

\begin{itemize}
  \item points fixes et stabilité de la dynamique moyenne du capital ;
  \item rôle de $\gamma$, $\delta$, $\sigma$, $K_0$, $k$ et $\lambda$ dans les
    échelles et nombres sans dimension ;
  \item conditions attendues de stationnarité, extinction ou croissance ;
  \item effets de la concavité sur les distributions de capital et de revenu ;
  \item articulation entre exposition au crédit, défaut et propagation des
    avalanches ;
  \item prédictions falsifiables permettant de distinguer une bifurcation, un
    transitoire long et une censure par garde-fou.
\end{itemize}

Le rapport final de sensibilité doit terminer par cette feuille de route et
indiquer quels résultats empiriques de M4B contraignent déjà l'analyse future.

\section{Vérification et conduite autonome}

Délègue les sous-tâches réellement indépendantes à des sous-agents, par exemple
l'inventaire des runs, la vérification statistique et l'audit des rapports.
Garde la conception expérimentale et la synthèse scientifique sous ta
responsabilité. À la fin de chaque phase majeure, fais vérifier par un
sous-agent à contexte frais : conformité au protocole, séparation exploration
et confirmation, calcul des métriques, diagnostic de \code{pop\_max},
provenance des figures et adéquation des conclusions aux données.

Avant chaque compte rendu de progression, vérifie chaque affirmation dans un
résultat d'outil de cette session. Si un test échoue, si un run est incomplet ou
si une étape est omise, dis-le explicitement. Ne présente jamais une intention
comme un travail accompli.

Travaille de façon autonome pour toutes les actions réversibles dans ce
périmètre. Ne demande l'utilisateur que si une action destructive ou
irréversible, un changement du moteur M4B, une extension de périmètre ou une
information que lui seul possède devient indispensable. Ne supprime aucun
résultat existant et ne modifie pas le moteur de référence.

\section{Définition de « terminé »}

\begin{briefbox}
La mission est terminée seulement lorsque les runs retenus sont exécutés ou
explicitement comptabilisés comme échecs ou censures, les contrôles sont
passés, les résultats sont reproductibles, le statut des franchissements de
\code{pop\_max} est établi, les figures sont sourcées et les trois rapports
LaTeX sont compilés et relus. Le rapport final doit conclure sur la sensibilité
de M4B, proposer prudemment les réductions éventuelles et présenter la feuille
de route M5 comme une perspective hors périmètre.
\end{briefbox}

Dans ton message final, commence par le résultat scientifique principal, puis
donne les chemins des rapports PDF et leurs limites éventuelles.

\end{document}
